% Figure IV.new --- keep-gain vs claw-back: the same panel, two critical
% rates. ch6-32 (register: should).
%
% Reader question: does it matter which slashing convention you use when
% quoting the audit rate needed for deterrence?
% Claim: at the same G=10, d=0.8, B=50, keep-gain needs rho*=0.25 but
% claw-back needs only rho*_c=0.2083 -- a platform that mismatches
% convention and formula is over- or under-auditing by about 17%.
% Evidence: rho*=G/(dB)=0.25; rho*_c=G/(d(G+B))=0.2083 [harbor3].
% Must distinguish: which bar belongs to which slashing convention, not
% just the two bare numbers.
% Grammar chosen: paired bar, one shared axis, both values computed from
% the same G, d, B (SKILL.md "quantity against a parameter", two-point
% case). Grammar rejected: a single number with a percent-difference
% callout -- it would not let a reader see both critical rates as
% positions on one common scale (Cleveland-McGill).
% Five-second test: a reader can see the claw-back bar is visibly shorter
% and name the roughly 17% gap.
\begin{figure}[H]
\centering
\pdfigurehue{pdgold}%
\begin{tikzpicture}[pd figure]
\begin{axis}[pd axis,
  width=6.0cm,height=3.6cm,
  ybar,bar width=1.1cm,
  symbolic x coords={keep-gain,claw-back},
  enlarge x limits=0.35,
  xtick=data,
  ymin=0,ymax=0.32,
  ytick={0,0.1,0.2,0.3},
  ylabel={critical rate $\rho^\star$},
  nodes near coords,
  every node near coord/.append style={pd label,anchor=south,yshift=2pt}]
  \addplot[pd focus fill] coordinates {(keep-gain,0.25) (claw-back,0.2083)};
\end{axis}
\node[pd note,anchor=north west,text width=8.0cm] at (-0.30,-2.55)
  {same $G=10$, $d=0.8$, $B=50$ in both bars; only the forfeiture on a
   caught cheat changes, $B$ versus $G+B$};
\end{tikzpicture}
\caption{\textbf{Quoting the wrong convention misprices deterrence by about
$17\%$.} Keep-gain forfeits only the bond $B$ on a caught cheat: critical
rate $\rho^\star=G/(dB)=0.25$. Claw-back also confiscates the gain $G$,
dropping the critical rate to $G/(d(G+B))=0.2083$. A platform that claws
back but budgets audits at the keep-gain rate over-audits; one that does not
claw back but budgets at the claw-back rate under-deters.
[internal; harbor3]}
\label{fig:he-deterrence-convention}
\end{figure}
